\chapter{Finding VI --- perception can be learned without a pretrained VAE: V-JEPA self-distillation}
\label{ch:jepa}

Chapters~\ref{ch:perception} and \ref{ch:recommendation} establish two facts that
together harden into a recommendation: on the short seven-step task perception cannot be
learned from reward (a from-scratch VAE collapses to always-press, orientation
$\Rsq=0.24$ against the pretrained VAE's $0.97$), and the cure is a \emph{frozen
pretrained VAE} front-end, declared one of two non-negotiables. This chapter is a third
perception result that qualifies that recommendation rather than overturning it. It shows
perception \emph{can} be learned jointly with reinforcement learning, with no pretrained or
frozen VAE, when the network is given a self-supervised auxiliary: a V-JEPA temporal
self-distillation loss (prototype-softmax, DINO/iBOT-style, EMA teacher, student$@t$
predicting teacher$@t{+}1$) trained jointly with the actor--critic. Paired with either time
to integrate (frame-repeat from raw pixels, $\sim0.90$ training rolling correct) or a
capable convolutional front-end (an SE-ResNet, $\sim0.87$ greedy correct with a sharp,
fully dissected covert-cueing algorithm), the auxiliary turns the perception bottleneck
into a solvable joint problem. The chapter also reports informative negatives that fence the
recipe: externalising a categorical softmax simplex as the working memory kills task
learning; a plain SmoothL1 latent-regression JEPA undergoes representational collapse; and
the matrix-affine two-LSTM feedback wiring collapses to chance under either the JEPA or a
VAE, isolating the fault to the wiring. The cross-cutting lesson is that the perceptual
front-end, not memory size, sets the attention strategy: a strong per-frame encoder yields a
sharp attention-driven solution, while learning from held pixels yields a diffuse,
recurrence-driven one. Because predicting future latents is arguably more biologically
plausible than supervised reconstruction-pretraining, this route is also the more neurally
defensible way to obtain a working front-end.

\section{The question Chapter 5 left open}
Chapter~\ref{ch:perception} draws a clean line between two questions that the earlier
architecture history had conflated: \emph{can the front-end see?} and \emph{does the
architecture work?} On the seven-step task the answer to the first was decisively negative
when perception was asked to emerge from reward alone. A VAE front-end trained end-to-end
from the reinforcement signal collapses to the \emph{always-press-at-the-deadline}
degenerate policy --- hit rate $1.000$, false-alarm rate $1.000$ on no-change trials,
overall accuracy $0.515$, reaction time pinned at frame $5$ --- and a linear probe confirms
the cause: orientation is decodable from the from-scratch encoder at only $\Rsq=0.24$ versus
$\Rsq=0.97$ for the pretrained VAE. Freezing a pretrained VAE rescues learning on the first
attempt (hit $0.79$, false-alarm $0.08$). Chapter~\ref{ch:recommendation} then elevates the
frozen pretrained VAE to one of \emph{two non-negotiables}.

That recommendation is correct \emph{as a statement about supervised reconstruction-
pretraining applied to a static, per-frame encoder}. It leaves a sharper question open: is
the pretrained VAE necessary because perception is genuinely unlearnable jointly with the
policy, or because reconstruction from reward is simply the wrong learning signal? This
chapter answers that perception is learnable jointly with the policy --- provided the
network carries a \emph{self-supervised predictive} auxiliary that does not depend on the
reward, and provided it is paired with one of two enablers that give the recurrence enough
perceptual purchase to use it. The result qualifies the non-negotiable: a frozen pretrained
VAE is one way to satisfy the perception requirement, not the only way.

\section{The auxiliary: V-JEPA temporal self-distillation}
The auxiliary is a video-JEPA self-distillation loss attached to the recurrent cell and
trained jointly with the actor--critic. A projection head reads the continuous xLSTM cell
output and emits, for each of $4$ heads, a softmax over $256$ prototypes. A teacher network
is an exponential moving average of the student (decay $0.996$); its outputs are centred
(centre momentum $0.9$) and sharpened (student temperature $0.1$; teacher temperature
warming from $0.04$ to $0.07$ over a $300$-iteration warm-up). The student's prototype
distribution at time $t$ is trained to predict the teacher's distribution at time $t{+}1$
--- a temporal, predict-the-next-latent target rather than a masked-spatial one. The loss
enters total optimisation only through a coefficient (\texttt{jepa-coef} $0.5$); crucially
it \emph{does not change the memory}, which stays a continuous xLSTM and is never itself
softmaxed. The prototype-softmax with centring and sharpening supplies the anti-collapse
machinery that, as Section~\ref{sec:jepa-negatives} shows, a naive latent loss lacks.

The diagnostic for whether the auxiliary is doing real work is the trajectory of the JEPA
loss. A loss that stays flat means the auxiliary never engages; a loss that drops to
$\approx0$ means representational collapse to a constant; a loss that descends to a healthy
non-zero plateau means structured, non-degenerate predictive learning. On the
frozen-colour-VAE control --- the exact paper network with the auxiliary bolted on --- the
auxiliary is a pure add-on and is not load-bearing: both feedback wirings reach $\sim0.85$
training rolling correct (peak $\sim0.90$) whether or not the JEPA engages, and \emph{whether
it engages at all is set by the feedback wiring}. Under cross-attention feedback the loss
falls $\sim3.7\times$ (first-200 mean $5.35\to$ late-200 mean $1.24$); under FiLM feedback it
sits inert ($5.32\to4.99$) --- though this FiLM run stopped early at $16{,}385$ of a
configured $50{,}000$ iterations, so its inertness is partly confounded with training length.
With perception already solved by the frozen VAE, the auxiliary neither helps nor hurts. Its
value appears only where there is no frozen VAE to solve perception for it.

Table~\ref{tab:jepa} summarises every variant studied in this chapter; the two routes that
make pixel-level learning work (Sections~\ref{sec:jepa-resultA} and \ref{sec:jepa-resultB})
and the three negatives that fence the recipe (Section~\ref{sec:jepa-negatives}) are
developed below.

\begin{table}[t]
\centering
\small
\begin{tabular}{@{}lllllrrll@{}}
\toprule
Model & Front-end & Feedback & JEPA flavour & Env & $d_{\text{mem}}$ & Iters & Rolling correct & Status \\
\midrule
paper\_jepa\_vae & frozen colour VAE & FiLM & proto-softmax & vda4 T7 & 1024 & 16385 & 0.846 (peak 0.90) & works (JEPA inert) \\
paper\_jepa\_vae & frozen colour VAE & crossattn1 & proto-softmax & vda4 T7 & 1024 & 49999 & 0.863 (peak 0.90) & works (JEPA engaged) \\
paper\_jepa\_fr\_novae & pixel (no VAE) & crossattn1 & proto-softmax & vda4 FR$\times$5 & 1024 & 33674 & 0.862 (final 0.895) & works \\
paper\_jepa\_fr\_novae & pixel (no VAE) & crossattn1 & proto-softmax & vda4 FR$\times$5 & 512 & 21240 & 0.807 (final 0.765) & works \\
paper\_jepa\_fr\_novae & pixel (no VAE) & crossattn1 & proto-softmax & vda4 FR$\times$5 & 256 & 15981 & 0.828 (final 0.848) & works \\
paper\_jepa\_fr\_novae & pixel (no VAE) & crossattn1 & proto-softmax & vda4 FR$\times$5 & 128 & 15327 & 0.828 (final 0.820) & works \\
paper\_jepa\_conv & SE-ResNet conv & crossattn1 & proto-softmax & vda4 T7 & 128 & 12989 & 0.870 (greedy 0.865) & works \\
paper\_jepa\_conv & SE-ResNet conv & affine\_ew & proto-softmax & vda4 T7 & 128 & 6662 & 0.867 (greedy $\sim$0.83) & works \\
paper\_jepa\_grid9 & SE-ResNet conv & crossattn1 & proto-softmax & vda9 3$\times$3 & 128 & $\sim$361 & 0.65 (final 0.695) & mid-training \\
paper\_jepa\_grid9 & SE-ResNet conv & affine\_ew & proto-softmax & vda9 3$\times$3 & 128 & $\sim$353 & 0.74 (final 0.803) & mid-training \\
paper\_jepa\_smem & pixel (no VAE) & crossattn1 & softmax-as-memory & vda4 FR$\times$5 & 1024 & 4211 & 0.003 (collapse) & failed \\
affine2lstm\_jepa & conv 4$\times$4, no VAE & matrix-affine & SmoothL1 latent & vda T7 & H2 cell & 18749 & 0.498 (chance) & collapsed \\
paper\_affine2lstm\_vae & conv 4$\times$4, VAE & matrix-affine & none (VAE ctrl) & vda T7 & H2 cell & 12196 & 0.496 (chance) & collapsed \\
\bottomrule
\end{tabular}
\caption{V-JEPA self-distillation variants. Rolling correct is the late-window training
rolling-correct mean unless a greedy-eval number is noted; greedy figures are held-out
(paper\_jepa\_conv only). The grid9 rows are mid-training (a few hundred iterations) and
report the last-200 mean (final in parentheses). No greedy-eval artifact exists for any run
except paper\_jepa\_conv crossattn1.}
\label{tab:jepa}
\end{table}

\section{Result A: learning from raw pixels via frame-repeat}
\label{sec:jepa-resultA}
The first route to a working pixel front-end gives the recurrence \emph{time to integrate}.
The cross-attention paper network is trained from raw pixels with no pretrained or frozen
VAE --- a standard linear patch-embed, one token per $2\times2$ cell --- with each logical
frame repeated five times, so the seven logical steps become $35$ physical steps (cue at
logical step $1$, change at logical step $5$, i.e.\ physical step $25$). With the V-JEPA
auxiliary engaged, this is one of the very few configurations in the program that learns the
change-detection task directly from pixels.

At $d_{\text{mem}}=1024$ the network reaches $0.895$ final and $0.862$ late-200-mean training
rolling correct, with the JEPA loss descending from $4.563$ to $0.458$ --- engaged, and not
collapsed to zero. A $d_{\text{mem}}$ sweep (1024/512/256/128) confirms the auxiliary engages
at every width (end-loss $0.458$/$0.762$/$0.526$/$0.516$) and that accuracy is
\emph{non-monotonic} in memory size: by late-200 mean the order is $d1024\,(0.862)>d256\,
(0.828)\approx d128\,(0.828)>d512\,(0.807)$. The widest memory is best and $512$ worst, but
the trend is not ``bigger is better'' --- consistent with a low-bandwidth task and with the
shorter runs ($d256$/$d128$ at $\sim15$--$16$k iterations versus $d1024$ at $\sim34$k) being
the least converged. All figures here are \emph{training rolling-correct}; no held-out
greedy-eval artifact exists for these runs.

The contrast with Chapter~\ref{ch:perception} is the point. A from-scratch VAE collapses on
the \emph{short, non-frame-repeated} seven-step task; the same family learns from pixels once
the held frames give the recurrence physical integration time and the JEPA shapes the
recurrent representation. The mechanistic signature of this route is \emph{diffuse,
recurrence-centric attention} (Fig.~\ref{fig:jepa-fr}): the cross-attention weights are
near-uniform across image and memory keys, with the only consistent structure a transient
per-query cue read on the cue frame; there is no change-lock at any $d_{\text{mem}}$. Change
detection lives in the LSTM recurrence, not in the attention weights.

\begin{figure}[t]
\centering
\includegraphics[width=0.78\linewidth]{fr_attn_memory.png}
\caption{\textbf{The frame-repeat pixel route is recurrence-centric.} Memory-key
cross-attention across the $35$ physical frames for the no-VAE frame-repeat model. The
allocation is near-uniform; the only structure is a transient cue-frame read. There is no
change-lock at any memory width --- detection is carried by the recurrent cell, not the
attention weights.}
\label{fig:jepa-fr}
\end{figure}

\section{Result B: the convolutional front-end and its sharp covert-cueing algorithm}
\label{sec:jepa-resultB}
The second route gives the network a \emph{capable per-frame encoder} instead of integration
time. A small SE-ResNet convolutional front-end ($\sim594$K parameters, trained end-to-end,
no VAE checkpoint) processes each $25\times25\times3$ patch through a $3\times3$ stem and
three squeeze-excite residual stages (GroupNorm, SiLU, downsample $25\to13\to7$, channels
$32\to64\to128$), global-average-pools, and emits a $128$-d embedding; the token is
$[\text{emb}_{128}\,\|\,\text{pos}_4\,\|\,\text{temporal}_8]=140$-d. On the plain seven-step
task ($\text{frame\_repeat}=1$, $d_{\text{mem}}=128$) with cross-attention feedback, the JEPA
loss descends from $\sim5.07$ (first-20 mean) to $\sim1.99$ (last-200 mean), a healthy
non-zero plateau, and the network reaches $0.8425$ final / $0.8695$ last-200-mean training
rolling correct. A held-out greedy evaluation at iteration $9{,}599$ (not the final iteration
$12{,}989$) gives correct $0.865$, hit $0.76$, correct-rejection $0.96$, false-alarm
$0.044$, reaction time $5.1\pm0.3$ --- a conservative detector and the only JEPA variant with
a measured greedy number.

Because a strong per-frame encoder solves perception without frame-repeat, the network can
afford an \emph{attention-driven} strategy, and the full mechanism is legible.
\textbf{Psychometric cueing} (Fig.~\ref{fig:jepa-psych}): forced detection rises with change
magnitude and saturates by $44^\circ$; validly cued changes are detected at lower magnitude
than invalid (P(detect) at $16^\circ$ $=0.58$ valid vs $0.43$ invalid; at $28^\circ$ $=0.99$
vs $0.94$), a classic left-shift cueing benefit driven by the spatial cue, not the displayed
validity ring. \textbf{Frame-gating} (Fig.~\ref{fig:jepa-framegate}): cross-attention reads
the image essentially only on the cue frame ($t{=}1$, image-key attention near ceiling) and
runs from memory on every Gabor/change frame (image-key $\approx0$); the current image is
never lost because it enters the cell via the residual $X^{(t)}$. \textbf{Cue-frame read}:
the three non-cued queries place $\sim1.00$ of their attention on the cued image patch,
broadcasting the cue to all four tokens. \textbf{Cue-timed memory change-lock}
(Fig.~\ref{fig:jepa-lock}): valid changes lock at $t{=}5$ (cue@S1 $+0.20$, cue@S4 $+0.24$)
and strongly at $t{=}6$ ($+0.75$, $+0.74$); invalid changes show \emph{no} lock at $t{=}5$
($-0.03$, $-0.02$) and lock only at $t{=}6$ ($+0.70$, $+0.69$) --- the cue buys one frame of
attentional latency, the attentional analogue of the behavioural cueing benefit.
\textbf{Validity/value invariance}: attention allocation is invariant to cue validity
(proportion $0.25$--$1.0$) and cue value (colour); validity and value act downstream at the
decision, a deployment-versus-use dissociation. \textbf{Temporal decoding}
(Fig.~\ref{fig:jepa-decode}): a four-fold cross-validated logistic probe on $\text{vec}(H)$
decodes cue position and value perfectly ($1.00$) from $t{=}1$, while change presence/location
is at chance until the change, then rises to $0.88$/$0.97$ at $t{=}5$ and $0.98$/$0.99$ at
$t{=}6$ --- change is localised in the recurrence at $t{=}5$, one frame before the
invalid-condition attention lock at $t{=}6$, so ``recurrence computes, attention reflects.''
\textbf{Causal memory read} (Fig.~\ref{fig:jepa-causal}): additively suppressing the cued
memory token costs little (P(detect) at $16^\circ$ $0.57\to0.47$); suppressing all four costs
more at high magnitude ($0.98\to0.84$ at $28^\circ$) but does not abolish detection, because
the changed Gabor rides in on the residual $X^{(t)}$. The memory read refines the decision by
$\sim0.1$--$0.15$; it is not the primary change-signal carrier.

\begin{figure}[t]
\centering
\includegraphics[width=0.62\linewidth]{dd4_psychometric.png}
\caption{\textbf{Psychometric cueing benefit.} P(declare change) versus change magnitude for
validly and invalidly cued changes. Valid cueing left-shifts the curve (e.g.\ $0.58$ vs
$0.43$ at $16^\circ$), saturating by $44^\circ$ --- the behavioural signature of covert
spatial attention.}
\label{fig:jepa-psych}
\end{figure}

\begin{figure}[t]
\centering
\includegraphics[width=0.72\linewidth]{dd2_framegate.png}
\caption{\textbf{Frame-gated attention.} Summed image-key cross-attention by frame. The
network reads the image only on the cue frame ($t{=}1$) and runs from memory thereafter; the
live image is preserved on the residual.}
\label{fig:jepa-framegate}
\end{figure}

\begin{figure}[t]
\centering
\includegraphics[width=0.82\linewidth]{dd2_maps_memory.png}
\caption{\textbf{Cue-timed memory change-lock.} Memory-key attention maps across frames.
Validly cued changes lock onto the changed location one frame earlier ($t{=}5$) than invalid
($t{=}6$) --- the attentional analogue of the behavioural cueing benefit.}
\label{fig:jepa-lock}
\end{figure}

\begin{figure}[t]
\centering
\includegraphics[width=0.62\linewidth]{dd5_decoding.png}
\caption{\textbf{Temporal decoding of the recurrent state.} Cross-validated balanced accuracy
of a logistic probe on $\text{vec}(H)$. Cue position/value are decodable from $t{=}1$; change
presence/location become decodable at $t{=}5$, one frame before the invalid-condition
attention lock --- recurrence computes, attention reflects.}
\label{fig:jepa-decode}
\end{figure}

\begin{figure}[t]
\centering
\includegraphics[width=0.62\linewidth]{dd5_causal.png}
\caption{\textbf{Causal perturbation of the memory read.} Additively suppressing memory
tokens on validly cued trials lowers detection modestly but does not abolish it: the changed
Gabor enters on the residual, so the memory read refines rather than carries the change
signal.}
\label{fig:jepa-causal}
\end{figure}

An element-wise-affine feedback variant (modulation $X'=\gamma(H)\odot X+\beta(H)$ then
self-attention over the four image patches, no separate memory-key column) reaches comparable
accuracy ($0.8625$ final / $0.8669$ last-200-mean rolling; greedy $\sim0.83$ at an early
iteration $4{,}199$) but solves the task differently: its change-lock is \emph{bottom-up},
locking onto the changed patch on the change frame for both valid ($1.00$) and invalid
($0.92$) changes --- any large orientation step pulls attention --- which predicts a smaller
cueing benefit. Its attention schedule is orient/release/re-acquire/lock (cued-patch attention
$0.73$ at cue, chance $0.22$ at blank, re-acquiring to $0.80$, saturating $1.0$ at a cued
change). It retains cue position and value at ceiling but \emph{perceives-then-discards} the
validity ring (decoded $0.99$ at $t{=}1$, collapsing to chance by Gabor onset $t{=}3$), the
mechanistic reason its behaviour cannot scale with validity. The two variants are opposite
poles of the same front-end: a memory-timed cueing lock versus a bottom-up salience grab.

\section{Informative negatives: what fences the recipe}
\label{sec:jepa-negatives}
Three negatives establish that the recipe is specific, not a generic ``add a JEPA'' fix.

\textbf{Softmax-as-memory fails.} Externalising the student's per-head softmax $S$
($4\times256=1024$-d) \emph{as} the recurrent memory --- so the same categorical simplex feeds
the cross-attention keys/values, the actor/critic readout, and the JEPA target --- kills task
learning on the no-VAE pixel frame-repeat task. The policy collapses: training rolling correct
peaks early at $0.453$ then decays monotonically to $0$ (late-200 mean $0.0029$). Critically,
the auxiliary itself trains fine (JEPA loss $4.55\to0.48$, a $\sim10\times$ drop, no collapse
to zero), so the failure is a clean dissociation: the representation learns while the policy
fails. A categorical probability simplex is too narrow a working-memory bottleneck; the
continuous cell must remain the working representation.

\textbf{Plain latent-regression JEPA collapses.} A simple SmoothL1 latent-regression auxiliary
(student projects $H2_t$ to predict EMA-teacher $H2_{t{+}1}$; no prototypes, no centring, no
sharpening) on a matrix-affine two-LSTM network exhibits the textbook latent collapse: the loss
falls from $0.489$ to $\sim0.11$ early in training, reaches a minimum of $0.00118$, and pins
near zero (late-2000 mean $0.0828$). With the predictor finding a near-constant target, the
auxiliary supplies no task signal and the policy sits at chance (late-200 mean rolling correct
$0.4975$, final $0.465$). This is the direct justification for the prototype-softmax/DINO-style
machinery in Results~A and~B: the anti-collapse apparatus is what makes the predictive
auxiliary useful.

\textbf{The collapse is the wiring, not the stabiliser.} The same matrix-affine two-LSTM
network collapses to always-wait chance under a \emph{VAE} as well (late-200 mean rolling
correct $0.4963$, final $0.45$). The JEPA did not rescue a wiring that fails regardless of
front-end stabiliser, so the fault is the matrix-affine two-LSTM feedback itself --- not
``affine'' in general, since the element-wise affine variant of Result~B reaches $\sim0.86$.
The negatives jointly localise the requirements: a continuous working memory, genuine
anti-collapse machinery in the auxiliary, and a feedback wiring that learns at all.

\section{Set-size-9 generalization (in progress)}
The conv-JEPA recipe is being carried from the four-stimulus $2\times2$ grid to a harder
nine-stimulus $3\times3$ grid (vda9: $75\times75$ image, coloured value cue, validity ring,
cue positions top-left/bottom-right). Both feedback wirings are above chance early and still
climbing. As of the current frontier (cross-attention at $\sim361$ iterations, element-wise
affine at $\sim353$), training rolling correct is $0.65$ (last-200 mean; $0.695$ final) for
cross-attention and $0.74$ (last-200 mean; $0.803$ final) for element-wise affine; the affine
variant showed an early JEPA descent (minimum $\sim2.46$ near iteration $94$) but the loss has
since risen back toward $\sim3.4$, so auxiliary engagement is not yet established for either
wiring at this set size. These are early, non-converged numbers --- a flat-looking curve this
early is not a ceiling --- so the result is reported as preliminary evidence that the no-VAE
conv-JEPA recipe transfers to the larger set size, not as a final accuracy.

\section{Cross-cutting lesson: architecture sets the attention strategy}
The two working routes reach comparable competence by opposite mechanisms, and the difference
is set by the \emph{perceptual front-end and the task temporal structure, not by memory size}.
A strong per-frame convolutional encoder (Result~B) lets the network afford an attention-driven
solution: a sharp, frame-gated algorithm that orients to the cue once, encodes it, monitors
from memory, and re-points to a change, with a cue-timed memory lock that yields a one-frame
attentional cueing benefit. Learning from held raw pixels (Result~A) instead yields a diffuse,
recurrence-centric solution in which the attention weights carry little structure and detection
lives in the LSTM. The $d_{\text{mem}}$ sweep confirms memory width is not the lever: accuracy
is non-monotonic in it. This sharpens the Chapter~\ref{ch:attention} reading that what carries
the computation --- attention versus recurrence --- is decided upstream of the memory.

For Chapter~\ref{ch:recommendation}, the consequence is a qualification, not a reversal. The
frozen pretrained VAE remains the simplest, most reliable, and best-characterised way to satisfy
the perception requirement, and remains the right default for the empirical manuscript. But the
requirement is now correctly stated as \emph{a front-end that reliably encodes the task
variable}, which can be met either by a frozen pretrained VAE \emph{or} by a self-supervised
predictive auxiliary (with proper anti-collapse machinery) paired with an enabler --- frame-
repeat or a capable conv encoder. The neuro relevance reinforces why this matters:
predict-the-future-latent self-supervision is arguably a more biologically plausible route to a
working visual front-end than supervised reconstruction-pretraining, so the JEPA route is not
merely an engineering alternative but the more defensible model of how a cortical front-end might
come to see.
